% Options for packages loaded elsewhere
\PassOptionsToPackage{unicode}{hyperref}
\PassOptionsToPackage{hyphens}{url}
%
\documentclass[
]{article}
\usepackage{amsmath,amssymb}
\usepackage{booktabs}
\usepackage{lmodern}
\usepackage{iftex}
\ifPDFTeX
  \usepackage[T1]{fontenc}
  \usepackage[utf8]{inputenc}
  \usepackage{textcomp} % provide euro and other symbols
\else % if luatex or xetex
  \usepackage{unicode-math}
  \defaultfontfeatures{Scale=MatchLowercase}
  \defaultfontfeatures[\rmfamily]{Ligatures=TeX,Scale=1}
\fi
% Use upquote if available, for straight quotes in verbatim environments
\IfFileExists{upquote.sty}{\usepackage{upquote}}{}
\IfFileExists{microtype.sty}{% use microtype if available
  \usepackage[]{microtype}
  \UseMicrotypeSet[protrusion]{basicmath} % disable protrusion for tt fonts
}{}
\makeatletter
\@ifundefined{KOMAClassName}{% if non-KOMA class
  \IfFileExists{parskip.sty}{%
    \usepackage{parskip}
  }{% else
    \setlength{\parindent}{0pt}
    \setlength{\parskip}{6pt plus 2pt minus 1pt}}
}{% if KOMA class
  \KOMAoptions{parskip=half}}
\makeatother
\usepackage{xcolor}
\IfFileExists{xurl.sty}{\usepackage{xurl}}{} % add URL line breaks if available
\IfFileExists{bookmark.sty}{\usepackage{bookmark}}{\usepackage{hyperref}}
\hypersetup{
  pdftitle={Effect of SAE on SUA - Main Results},
  hidelinks,
  pdfcreator={LaTeX via pandoc}}
\urlstyle{same} % disable monospaced font for URLs
\usepackage[margin=1in]{geometry}
\usepackage{graphicx}
\makeatletter
\def\maxwidth{\ifdim\Gin@nat@width>\linewidth\linewidth\else\Gin@nat@width\fi}
\def\maxheight{\ifdim\Gin@nat@height>\textheight\textheight\else\Gin@nat@height\fi}
\makeatother
% Scale images if necessary, so that they will not overflow the page
% margins by default, and it is still possible to overwrite the defaults
% using explicit options in \includegraphics[width, height, ...]{}
\setkeys{Gin}{width=\maxwidth,height=\maxheight,keepaspectratio}
% Set default figure placement to htbp
\makeatletter
\def\fps@figure{htbp}
\makeatother
\setlength{\emergencystretch}{3em} % prevent overfull lines
\providecommand{\tightlist}{%
  \setlength{\itemsep}{0pt}\setlength{\parskip}{0pt}}
\setcounter{secnumdepth}{-\maxdimen} % remove section numbering
\ifLuaTeX
  \usepackage{selnolig}  % disable illegal ligatures
\fi

\title{Effect of SAE on SUA - Main Results}
\author{}
\date{\vspace{-2.5em}}

\begin{document}
\maketitle

\hypertarget{data-and-treatment}{%
\subsection{Data and treatment}\label{data-and-treatment}}

\hypertarget{headline-table}{%
\subsection{Headline table}\label{headline-table}}

Graduated high school on time, having spent all four years (grades 9 to
12) in the same school. One observation per school x 9th-grade cohort;
each outcome is the average over the entire graduating class.

\begin{table}[htbp]
   \caption{ Regression Results}\label{tab: main results}
   \centerline{\scalebox{0.85}{\begin{tabular}{lcccccc}
      \toprule
                   & GPA & P(Take Exam) & Average LM & Applied & Assigned & Enrolled\\
                   & (1) & (2) & (3) & (4) & (5) & (6)\\
      \midrule
      Treated & -0.005 & -0.026$^{**}$ & -0.016$^{***}$ & 0.045$^{***}$ & 0.051$^{***}$ & 0.041$^{***}$\\
       & (0.018) & (0.011) & (0.004) & (0.006) & (0.005) & (0.009)\\
      \midrule
      Mean (pre-treatment) & 5.687 & 0.832 & 0.413 & 0.432 & 0.328 & 0.265\\
      Observations & 28,483 & 28,483 & 28,483 & 28,483 & 28,483 & 28,483\\
      \bottomrule
   \end{tabular}}}
\end{table}

\hypertarget{robustness-relaxed-sample}{%
\subsection{Robustness: relaxed
sample}\label{robustness-relaxed-sample}}

Same design, relaxing the headline sample in two ways. First, students
no longer have to graduate on time: anyone who graduates counts, keyed
to their graduating class. Second, the four-years-in-the-same-school
requirement is relaxed to at most one publicly funded school in high
school (so international students and private-to-public movers enter;
public-to-public switchers drop, since there is no single class to
attribute them to). Students entering 9th grade before their school
joined SAE do not count as treated.

\begin{table}[htbp]
   \caption{Regression Results: Relaxed Sample}\label{tab: robustness relaxed sample}
   \centerline{\scalebox{0.85}{\begin{tabular}{lcccccc}
      \toprule
                   & GPA & P(Take Exam) & Average LM & Applied & Assigned & Enrolled\\
                   & (1) & (2) & (3) & (4) & (5) & (6)\\
      \midrule
      Treated & -0.016 & -0.020 & -0.013$^{***}$ & 0.048$^{***}$ & 0.054$^{***}$ & 0.040$^{***}$\\
       & (0.028) & (0.014) & (0.003) & (0.006) & (0.005) & (0.012)\\
      \midrule
      Mean (pre-treatment) & 5.644 & 0.821 & 0.405 & 0.421 & 0.318 & 0.257\\
      Observations & 28,688 & 28,688 & 28,688 & 28,688 & 28,688 & 28,688\\
      \bottomrule
   \end{tabular}}}
\end{table}

\hypertarget{heterogeneity-gender}{%
\subsection{Heterogeneity: gender}\label{heterogeneity-gender}}

Headline sample, class means computed separately over the female and
male students of each class (students with unknown sex are dropped); the
control is the subgroup class size.

\begin{table}[htbp]
   \caption{Regression Results by Gender}\label{tab: results by gender}
   \centerline{\scalebox{0.85}{\begin{tabular}{lcccccccccccc}
      \toprule
      & \multicolumn{2}{c}{GPA}
      & \multicolumn{2}{c}{P(Take Exam)}
      & \multicolumn{2}{c}{Average LM}
      & \multicolumn{2}{c}{Applied}
      & \multicolumn{2}{c}{Assigned}
      & \multicolumn{2}{c}{Enrolled} \\
      \cmidrule(lr){2-3}\cmidrule(lr){4-5}\cmidrule(lr){6-7}\cmidrule(lr){8-9}\cmidrule(lr){10-11}\cmidrule(lr){12-13}
                   & Female & Male & Female & Male & Female & Male & Female & Male & Female & Male & Female & Male \\
                   & (1) & (2) & (3) & (4) & (5) & (6) & (7) & (8) & (9) & (10) & (11) & (12)\\
      \midrule
      Treated & -0.005 & -0.009 & -0.013 & -0.030$^{*}$ & -0.011$^{***}$ & -0.015$^{***}$ & 0.078$^{***}$ & 0.020$^{**}$ & 0.075$^{***}$ & 0.030$^{***}$ & 0.064$^{***}$ & 0.024$^{***}$\\
       & (0.013) & (0.013) & (0.014) & (0.018) & (0.004) & (0.005) & (0.008) & (0.009) & (0.007) & (0.006) & (0.017) & (0.008)\\
      \midrule
      Mean (pre-treatment) & 5.756 & 5.602 & 0.852 & 0.810 & 0.410 & 0.415 & 0.464 & 0.396 & 0.335 & 0.319 & 0.265 & 0.265\\
      Observations & 27,708 & 27,240 & 27,708 & 27,240 & 27,708 & 27,240 & 27,708 & 27,240 & 27,708 & 27,240 & 27,708 & 27,240\\
      \bottomrule
   \end{tabular}}}
\end{table}

\hypertarget{heterogeneity-school-vulnerability}{%
\subsection{Heterogeneity: school
vulnerability}\label{heterogeneity-school-vulnerability}}

Treated schools split at the within-region median of the 2015 IVE-SINAE
vulnerability index, taken over treated schools (treated schools with no
index are dropped); never-treated private schools serve as the
comparison group in both arms.

\begin{table}[htbp]
   \caption{Regression Results by School Vulnerability}\label{tab: results by vulnerability}
   \centerline{\scalebox{0.85}{\begin{tabular}{lcccccccccccc}
      \toprule
      & \multicolumn{2}{c}{GPA}
      & \multicolumn{2}{c}{P(Take Exam)}
      & \multicolumn{2}{c}{Average LM}
      & \multicolumn{2}{c}{Applied}
      & \multicolumn{2}{c}{Assigned}
      & \multicolumn{2}{c}{Enrolled} \\
      \cmidrule(lr){2-3}\cmidrule(lr){4-5}\cmidrule(lr){6-7}\cmidrule(lr){8-9}\cmidrule(lr){10-11}\cmidrule(lr){12-13}
                   & Low & High & Low & High & Low & High & Low & High & Low & High & Low & High \\
                   & (1) & (2) & (3) & (4) & (5) & (6) & (7) & (8) & (9) & (10) & (11) & (12)\\
      \midrule
      Treated & 0.013 & -0.030 & -0.025$^{***}$ & -0.026 & -0.039$^{***}$ & 0.008 & 0.025$^{***}$ & 0.068$^{***}$ & 0.027$^{***}$ & 0.077$^{***}$ & 0.015$^{**}$ & 0.068$^{***}$\\
       & (0.010) & (0.032) & (0.005) & (0.019) & (0.004) & (0.005) & (0.004) & (0.008) & (0.005) & (0.007) & (0.006) & (0.010)\\
      \midrule
      Mean (pre-treatment) & 5.796 & 5.569 & 0.939 & 0.725 & 0.561 & 0.265 & 0.589 & 0.273 & 0.467 & 0.187 & 0.379 & 0.150\\
      Observations & 15,633 & 15,504 & 15,633 & 15,504 & 15,633 & 15,504 & 15,633 & 15,504 & 15,633 & 15,504 & 15,633 & 15,504\\
      \bottomrule
   \end{tabular}}}
\end{table}

\hypertarget{heterogeneity-program-selectivity}{%
\subsection{Heterogeneity: program
selectivity}\label{heterogeneity-program-selectivity}}

Share of the graduating class assigned to a program, by the selectivity
of the program (programs that fill their seats, and cutoff-decile bins);
each flag is zero for students not assigned, so the shares are over the
entire class.

\begin{table}[htbp]
   \caption{Regression Results by Program Selectivity}\label{tab: results by program selectivity}
   \centerline{\scalebox{0.85}{\begin{tabular}{lcccccccc}
      \toprule
                   & All & Full & Not Full & Top 10 & Top 25 & Top 50 & Top 75 & Top 90\\
                   & (1) & (2) & (3) & (4) & (5) & (6) & (7) & (8)\\
      \midrule
      Treated & 0.051$^{***}$ & -0.003 & 0.055$^{***}$ & -0.030$^{***}$ & -0.004 & 0.036$^{*}$ & 0.051$^{***}$ & 0.051$^{***}$\\
       & (0.005) & (0.009) & (0.007) & (0.011) & (0.011) & (0.019) & (0.005) & (0.005)\\
      \midrule
      Mean (pre-treatment) & 0.328 & 0.250 & 0.078 & 0.047 & 0.115 & 0.250 & 0.322 & 0.328\\
      Observations & 28,483 & 28,483 & 28,483 & 28,483 & 28,483 & 28,483 & 28,483 & 28,483\\
      \bottomrule
   \end{tabular}}}
\end{table}

\end{document}
